% Options for packages loaded elsewhere
% Options for packages loaded elsewhere
\PassOptionsToPackage{unicode}{hyperref}
\PassOptionsToPackage{hyphens}{url}
\PassOptionsToPackage{dvipsnames,svgnames,x11names}{xcolor}
%
\documentclass[
  a4paper,
]{article}
\usepackage{xcolor}
\usepackage[top=1in,bottom=1in,left=0.9in,right=0.9in]{geometry}
\usepackage{amsmath,amssymb}
\setcounter{secnumdepth}{3}
\usepackage{iftex}
\ifPDFTeX
  \usepackage[T1]{fontenc}
  \usepackage[utf8]{inputenc}
  \usepackage{textcomp} % provide euro and other symbols
\else % if luatex or xetex
  \usepackage{unicode-math} % this also loads fontspec
  \defaultfontfeatures{Scale=MatchLowercase}
  \defaultfontfeatures[\rmfamily]{Ligatures=TeX,Scale=1}
\fi
\usepackage{lmodern}
\ifPDFTeX\else
  % xetex/luatex font selection
\fi
% Use upquote if available, for straight quotes in verbatim environments
\IfFileExists{upquote.sty}{\usepackage{upquote}}{}
\IfFileExists{microtype.sty}{% use microtype if available
  \usepackage[]{microtype}
  \UseMicrotypeSet[protrusion]{basicmath} % disable protrusion for tt fonts
}{}
\makeatletter
\@ifundefined{KOMAClassName}{% if non-KOMA class
  \IfFileExists{parskip.sty}{%
    \usepackage{parskip}
  }{% else
    \setlength{\parindent}{0pt}
    \setlength{\parskip}{6pt plus 2pt minus 1pt}}
}{% if KOMA class
  \KOMAoptions{parskip=half}}
\makeatother
% Make \paragraph and \subparagraph free-standing
\makeatletter
\ifx\paragraph\undefined\else
  \let\oldparagraph\paragraph
  \renewcommand{\paragraph}{
    \@ifstar
      \xxxParagraphStar
      \xxxParagraphNoStar
  }
  \newcommand{\xxxParagraphStar}[1]{\oldparagraph*{#1}\mbox{}}
  \newcommand{\xxxParagraphNoStar}[1]{\oldparagraph{#1}\mbox{}}
\fi
\ifx\subparagraph\undefined\else
  \let\oldsubparagraph\subparagraph
  \renewcommand{\subparagraph}{
    \@ifstar
      \xxxSubParagraphStar
      \xxxSubParagraphNoStar
  }
  \newcommand{\xxxSubParagraphStar}[1]{\oldsubparagraph*{#1}\mbox{}}
  \newcommand{\xxxSubParagraphNoStar}[1]{\oldsubparagraph{#1}\mbox{}}
\fi
\makeatother


\usepackage{longtable,booktabs,array}
\newcounter{none} % for unnumbered tables
\usepackage{calc} % for calculating minipage widths
% Correct order of tables after \paragraph or \subparagraph
\usepackage{etoolbox}
\makeatletter
\patchcmd\longtable{\par}{\if@noskipsec\mbox{}\fi\par}{}{}
\makeatother
% Allow footnotes in longtable head/foot
\IfFileExists{footnotehyper.sty}{\usepackage{footnotehyper}}{\usepackage{footnote}}
\makesavenoteenv{longtable}
\usepackage{graphicx}
\makeatletter
\newsavebox\pandoc@box
\newcommand*\pandocbounded[1]{% scales image to fit in text height/width
  \sbox\pandoc@box{#1}%
  \Gscale@div\@tempa{\textheight}{\dimexpr\ht\pandoc@box+\dp\pandoc@box\relax}%
  \Gscale@div\@tempb{\linewidth}{\wd\pandoc@box}%
  \ifdim\@tempb\p@<\@tempa\p@\let\@tempa\@tempb\fi% select the smaller of both
  \ifdim\@tempa\p@<\p@\scalebox{\@tempa}{\usebox\pandoc@box}%
  \else\usebox{\pandoc@box}%
  \fi%
}
% Set default figure placement to htbp
\def\fps@figure{htbp}
\makeatother





\setlength{\emergencystretch}{3em} % prevent overfull lines

\providecommand{\tightlist}{%
  \setlength{\itemsep}{0pt}\setlength{\parskip}{0pt}}



 


\usepackage{pdflscape}
\hyphenpenalty=10000
\exhyphenpenalty=10000
\emergencystretch=3em
\makeatletter
\renewcommand*\l@subsection{\@dottedtocline{2}{1.5em}{3.8em}}
\renewcommand*\l@subsubsection{\@dottedtocline{3}{5.3em}{4.8em}}
\makeatother
\makeatletter
\@ifpackageloaded{caption}{}{\usepackage{caption}}
\AtBeginDocument{%
\ifdefined\contentsname
  \renewcommand*\contentsname{Table of contents}
\else
  \newcommand\contentsname{Table of contents}
\fi
\ifdefined\listfigurename
  \renewcommand*\listfigurename{List of Figures}
\else
  \newcommand\listfigurename{List of Figures}
\fi
\ifdefined\listtablename
  \renewcommand*\listtablename{List of Tables}
\else
  \newcommand\listtablename{List of Tables}
\fi
\ifdefined\figurename
  \renewcommand*\figurename{Figure}
\else
  \newcommand\figurename{Figure}
\fi
\ifdefined\tablename
  \renewcommand*\tablename{Table}
\else
  \newcommand\tablename{Table}
\fi
}
\@ifpackageloaded{float}{}{\usepackage{float}}
\floatstyle{ruled}
\@ifundefined{c@chapter}{\newfloat{codelisting}{h}{lop}}{\newfloat{codelisting}{h}{lop}[chapter]}
\floatname{codelisting}{Listing}
\newcommand*\listoflistings{\listof{codelisting}{List of Listings}}
\makeatother
\makeatletter
\makeatother
\makeatletter
\@ifpackageloaded{caption}{}{\usepackage{caption}}
\@ifpackageloaded{subcaption}{}{\usepackage{subcaption}}
\makeatother
\usepackage{bookmark}
\IfFileExists{xurl.sty}{\usepackage{xurl}}{} % add URL line breaks if available
\urlstyle{same}
\hypersetup{
  pdftitle={Python for Applied Real Estate and Climate-Risk Analytics},
  pdfauthor={James Daniel (PhD)},
  colorlinks=true,
  linkcolor={blue},
  filecolor={Maroon},
  citecolor={Blue},
  urlcolor={Blue},
  pdfcreator={LaTeX via pandoc}}


\title{Python for Applied Real Estate and Climate-Risk Analytics}
\usepackage{etoolbox}
\makeatletter
\providecommand{\subtitle}[1]{% add subtitle to \maketitle
  \apptocmd{\@title}{\par {\large #1 \par}}{}{}
}
\makeatother
\subtitle{A 12-Week Professional Curriculum for Flood Risk, Physical
Climate Risk, Transition Risk, Valuation, REITs and Investment
Decision-Making}
\author{James Daniel (PhD)}
\date{June 13, 2026}
\begin{document}
\maketitle

\renewcommand*\contentsname{Table of contents}
{
\setcounter{tocdepth}{2}
\tableofcontents
}

\section{Programme Title}\label{programme-title}

\textbf{Python for Applied Real Estate and Climate-Risk Analytics: Flood
Risk, Physical Risk, Transition Risk, Valuation and Investment
Decision-Making}

This professional curriculum is a 12-week / 90-day applied Python
programme for real estate finance, climate-risk analytics, investment
research and professional decision-support. It is designed for a serious
academic and professional learner who needs Python not as a generic
programming language, but as a practical research and analytics
environment for collecting, cleaning, modelling, visualising and
communicating climate-risk evidence in real estate markets.

\section{Programme Rationale}\label{programme-rationale}

Real estate markets are increasingly exposed to physical climate risk,
flood risk, transition risk, insurance stress, mortgageability
constraints, infrastructure fragility, energy-performance regulation,
ESG reporting expectations and climate-adjusted investment
decision-making. These pressures are no longer peripheral to valuation,
finance, portfolio strategy or policy analysis. They influence pricing,
liquidity, asset resilience, investment risk, stranded-asset exposure,
REIT volatility and the long-term credibility of real estate market
evidence.

For a PhD-trained real estate investment and finance professional, the
central challenge is not simply learning Python syntax. The challenge is
converting advanced conceptual, quantitative and domain knowledge into a
reproducible computational workflow. This includes the ability to
assemble property, transaction, climate, flood-risk, geospatial,
financial and textual datasets; clean and integrate heterogeneous data
sources; build defensible analytical models; produce professional
visualisations and dashboards; communicate results to academic,
investor, policy, valuation and consultancy audiences; and maintain
reproducible evidence pipelines suitable for journal articles,
grant-funded projects, consultancy reports and future analytic products.

This curriculum therefore rejects a generic ``learn Python'' structure.
It begins with foundations because the learner's programming experience
is beginner-to-developing, but it moves quickly into applied analytics
because the learner's domain reasoning is already advanced. The
programme uses real estate and climate-risk problems from the beginning,
so that variables, functions, loops, data frames, regression models,
machine-learning workflows, geospatial joins, time-series analysis, NLP
and dashboards are introduced through meaningful professional problems
rather than abstract coding drills.

\section{Learner Profile}\label{learner-profile}

The intended learner is a PhD-trained academic and professional in real
estate investment and finance, with expertise in property market
analysis, empirical research design, regression-based analysis,
financial and economic interpretation, valuation, investment
decision-making and climate-risk research.

The learner's current professional and research interests include:

\begin{itemize}
\tightlist
\item
  flood-risk capitalisation in housing markets;
\item
  physical climate risk and property valuation;
\item
  transition risk, stranded assets and ESG regulation;
\item
  real estate pricing, valuation and investment analytics;
\item
  infrastructure resilience and urban resilience;
\item
  housing markets, REITs and climate-adjusted investment strategies;
\item
  insurance affordability, mortgageability and property market
  liquidity;
\item
  dashboards for investors, valuers, insurers, policymakers, local
  authorities and asset managers.
\end{itemize}

The learner has advanced conceptual, statistical, financial and research
knowledge, but should be treated as beginner-to-developing in Python
programming. The curriculum must therefore provide clear scaffolding for
syntax and computational thinking while respecting the learner's ability
to engage with advanced research questions, modelling logic and
interpretation.

\section{Programme Aim}\label{programme-aim}

The aim of this programme is to develop the learner's capacity to use
Python confidently for applied real estate and climate-risk research,
professional analytics, consultancy, grant-funded projects, journal
publications and future dashboard or product development.

By the end of the programme, the learner should be able to move from a
conceptually formulated real estate climate-risk problem to a complete
analytical workflow: identifying suitable data, importing and cleaning
datasets, integrating property and climate-risk indicators, conducting
exploratory analysis, building statistical and machine-learning models,
analysing time-series and geospatial patterns, extracting
transition-risk signals from text, building professional visualisations
and dashboards, documenting the workflow through Quarto and GitHub, and
communicating findings in a rigorous, reproducible and decision-relevant
format.

\section{Learning Outcomes}\label{learning-outcomes}

By the end of the 12-week programme, the learner should be able to:

\begin{enumerate}
\def\labelenumi{\arabic{enumi}.}
\tightlist
\item
  Set up and manage a professional Python analytics environment using
  Python, JupyterLab, VS Code, Quarto, Git and GitHub.
\item
  Apply Python foundations, including variables, data types, functions,
  loops, conditionals and basic control flow, to real estate and
  climate-risk examples.
\item
  Import, clean, validate and manage Excel, CSV, API, web-sourced,
  geospatial, financial and textual datasets.
\item
  Use \texttt{pandas} and \texttt{numpy} to structure, transform, merge
  and summarise property, transaction, valuation, climate-risk and
  financial datasets.
\item
  Conduct exploratory data analysis on real estate and climate-risk
  data, including missing-data assessment, distributional analysis,
  group comparisons, correlation patterns and outlier detection.
\item
  Produce professional charts, tables, maps and interactive
  visualisations using \texttt{matplotlib}, \texttt{seaborn},
  \texttt{plotly}, \texttt{geopandas} and \texttt{folium}.
\item
  Estimate and interpret statistical and econometric models relevant to
  real estate pricing, flood-risk capitalisation, valuation adjustment
  and climate-risk exposure.
\item
  Apply machine-learning models for prediction and classification tasks
  such as exposure scoring, price-discount estimation, vulnerability
  classification, investment-risk assessment and asset-level resilience
  screening.
\item
  Analyse time-series data for housing prices, REIT returns, interest
  rates, inflation, volatility, climate events, infrastructure shocks
  and energy-market shocks.
\item
  Conduct geospatial analysis using property locations, flood-risk
  layers, local authority boundaries, vulnerability indicators and
  climate-exposure measures.
\item
  Use introductory natural language processing to analyse ESG
  disclosures, planning reports, policy documents and market sentiment
  for transition-risk signals.
\item
  Build a prototype climate-risk analytics dashboard using Streamlit or
  Dash.
\item
  Apply reproducible workflow principles, including project folder
  structure, notebook discipline, Quarto reporting, Git/GitHub version
  control, documentation and environment management.
\item
  Communicate real estate climate-risk findings to academic, investor,
  policy, valuation and consultancy audiences through technical reports,
  dashboards and professional presentations.
\item
  Design a progression route from beginner-to-intermediate Python
  competence toward advanced econometrics, spatial modelling, causal
  inference, climate scenario analysis, portfolio analytics and research
  publication.
\end{enumerate}

\section{Curriculum Design
Principles}\label{curriculum-design-principles}

The curriculum is guided by eight principles.

First, \textbf{domain specificity} is treated as non-negotiable. Python
is taught through real estate, climate-risk and investment problems from
the first week.

Second, \textbf{progressive complexity} structures the programme. The
learner starts with syntax and computational logic, then moves into data
wrangling, visualisation, regression, machine learning, time series,
geospatial analytics, NLP, dashboards and reproducible deployment.

Third, \textbf{research alignment} ensures that weekly activities
support the learner's academic and professional goals: journal articles,
grant-funded projects, consultancy outputs, investment reports and
future analytical products.

Fourth, \textbf{applied learning} dominates the design. Each week
produces a portfolio artifact rather than only theoretical knowledge.

Fifth, \textbf{reproducibility} is built from the beginning through
folder structures, notebooks, Quarto reports, version control,
documented assumptions and transparent workflows.

Sixth, \textbf{professional communication} is treated as part of
analytics. The learner must not only run code but also explain results
in ways that are meaningful for valuers, investors, policymakers,
insurers, local authorities, asset managers and academic reviewers.

Seventh, \textbf{project-based assessment} ensures that evidence of
competence comes from authentic tasks: cleaning datasets, estimating
models, mapping exposure, building dashboards and presenting defensible
findings.

Eighth, \textbf{portfolio development} ensures that by the end of the
programme the learner has a coherent body of work that can support
research, consultancy, teaching, grant applications and product
development.

\section{Required Tools and Technical
Setup}\label{required-tools-and-technical-setup}

{\def\LTcaptype{none} % do not increment counter
\begin{longtable}[]{@{}
  >{\raggedright\arraybackslash}p{(\linewidth - 4\tabcolsep) * \real{0.3333}}
  >{\raggedright\arraybackslash}p{(\linewidth - 4\tabcolsep) * \real{0.3333}}
  >{\raggedright\arraybackslash}p{(\linewidth - 4\tabcolsep) * \real{0.3333}}@{}}
\toprule\noalign{}
\begin{minipage}[b]{\linewidth}\raggedright
Tool
\end{minipage} & \begin{minipage}[b]{\linewidth}\raggedright
Curriculum Use
\end{minipage} & \begin{minipage}[b]{\linewidth}\raggedright
Rationale
\end{minipage} \\
\midrule\noalign{}
\endhead
\bottomrule\noalign{}
\endlastfoot
Python 3.x & Core programming language & Provides the computational
foundation for data analysis, modelling, visualisation and
automation. \\
Anaconda or Miniconda & Environment and package management & Supports
installation of scientific Python libraries and reproducible
environments. Miniconda is preferable for lightweight, controlled
setups. \\
JupyterLab & Interactive coding and analysis & Suitable for exploratory
work, teaching, code demonstration, data cleaning and model
development. \\
VS Code & Professional code editing & Supports script development,
project navigation, Git integration and transition from notebooks to
production-oriented workflows. \\
Git & Version control & Enables tracking of changes, reproducibility,
rollback and disciplined project management. \\
GitHub & Remote repository and portfolio platform & Supports
professional documentation, project sharing, collaboration and
public/private portfolio development. \\
Quarto & Reproducible reporting & Allows code, explanation, tables,
figures and interpretation to be integrated into HTML/PDF reports. \\
Streamlit or Dash & Dashboard development & Enables rapid creation of
interactive tools for property, flood-risk, transition-risk and
investment decision-support. \\
Spreadsheet software & Data inspection and stakeholder compatibility &
Many property, valuation and consultancy datasets arrive in Excel; the
learner must understand how to move between spreadsheet and Python
workflows. \\
Browser and API tools & Data acquisition & Supports access to public
datasets, financial data, climate portals, policy documents and
web-sourced data where legally and ethically permitted. \\
Package-management file & Reproducibility & An \texttt{environment.yml},
\texttt{requirements.txt} or equivalent file should document
dependencies. \\
\end{longtable}
}

Recommended project structure:

{\def\LTcaptype{none} % do not increment counter
\begin{longtable}[]{@{}
  >{\raggedright\arraybackslash}p{(\linewidth - 2\tabcolsep) * \real{0.5000}}
  >{\raggedright\arraybackslash}p{(\linewidth - 2\tabcolsep) * \real{0.5000}}@{}}
\toprule\noalign{}
\begin{minipage}[b]{\linewidth}\raggedright
Folder / File
\end{minipage} & \begin{minipage}[b]{\linewidth}\raggedright
Purpose
\end{minipage} \\
\midrule\noalign{}
\endhead
\bottomrule\noalign{}
\endlastfoot
\texttt{data\_raw/} & Original unmodified data, subject to licensing and
confidentiality constraints. \\
\texttt{data\_processed/} & Cleaned and analysis-ready data. \\
\texttt{notebooks/} & Exploratory Jupyter notebooks. \\
\texttt{scripts/} & Reusable Python scripts for cleaning, modelling and
visualisation. \\
\texttt{reports/} & Quarto reports and exported outputs. \\
\texttt{dashboards/} & Streamlit or Dash application files. \\
\texttt{figures/} & Exported charts, maps and dashboard images. \\
\texttt{models/} & Saved model artifacts where appropriate. \\
\texttt{docs/} & Documentation, data dictionary, methodology notes and
assumptions. \\
\texttt{README.md} & Project overview and usage instructions. \\
\texttt{requirements.txt} or \texttt{environment.yml} & Package and
environment specification. \\
\end{longtable}
}

\section{Core Python Libraries}\label{core-python-libraries}

{\def\LTcaptype{none} % do not increment counter
\begin{longtable}[]{@{}
  >{\raggedright\arraybackslash}p{(\linewidth - 4\tabcolsep) * \real{0.3333}}
  >{\raggedright\arraybackslash}p{(\linewidth - 4\tabcolsep) * \real{0.3333}}
  >{\raggedright\arraybackslash}p{(\linewidth - 4\tabcolsep) * \real{0.3333}}@{}}
\toprule\noalign{}
\begin{minipage}[b]{\linewidth}\raggedright
Library / Tool
\end{minipage} & \begin{minipage}[b]{\linewidth}\raggedright
Curriculum Purpose
\end{minipage} & \begin{minipage}[b]{\linewidth}\raggedright
Applied Use in Real Estate and Climate-Risk Analytics
\end{minipage} \\
\midrule\noalign{}
\endhead
\bottomrule\noalign{}
\endlastfoot
\texttt{pandas} & Data manipulation and analysis & Cleaning transaction
records, merging property and flood-risk indicators, grouping by local
authority, handling valuation data. \\
\texttt{numpy} & Numerical computing & Arrays, mathematical operations,
transformations, simulation inputs and model preparation. \\
\texttt{matplotlib} & Static plotting & Publication-style charts for
house prices, REIT returns, exposure distributions and model
diagnostics. \\
\texttt{seaborn} & Statistical visualisation & Distribution plots,
correlation heatmaps, regression plots and comparative risk
visualisations. \\
\texttt{plotly} & Interactive visualisation & Interactive charts for
dashboards, investment summaries, time-series exploration and
stakeholder communication. \\
\texttt{scipy} & Scientific and statistical computation & Statistical
tests, distributions and numerical operations where required. \\
\texttt{statsmodels} & Statistical and econometric modelling & OLS
regression, model diagnostics, time-series models and interpretable real
estate pricing models. \\
\texttt{scikit-learn} & Machine learning & Prediction, classification,
model evaluation, pipelines and preprocessing for exposure,
vulnerability and price-discount models. \\
\texttt{geopandas} & Geospatial data analysis & Reading
shapefiles/GeoJSON, spatial joins, local authority boundaries and
flood-exposure mapping. \\
\texttt{shapely} & Geometric operations & Point-in-polygon analysis,
buffers, intersections and distance-based spatial features. \\
\texttt{folium} & Interactive web maps & Flood-risk maps, property
exposure maps and local authority risk dashboards. \\
\texttt{rasterio} & Raster geospatial data & Climate or flood raster
layers where available and legally accessible. \\
\texttt{xarray} & Multi-dimensional climate data & Climate grids,
time-indexed climate variables and scenario-oriented datasets where
appropriate. \\
\texttt{requests} & API access & Downloading public datasets, financial
data, policy data and open-data records where API access is
permitted. \\
\texttt{beautifulsoup4} & Web parsing & Extracting text or tables from
web pages where legally and ethically permitted. \\
\texttt{nltk} & Introductory NLP & Tokenisation, stopword removal and
basic text analysis for ESG and policy documents. \\
\texttt{spaCy} & Applied NLP & Named entity recognition, linguistic
preprocessing and document-level analysis. \\
\texttt{scikit-learn} text tools & Text vectorisation and classification
& TF-IDF, document classification and transition-risk signal
extraction. \\
\texttt{yfinance} or finance-data equivalent & Financial time-series
data & REIT prices, market indices and finance examples, subject to
verification and data-use terms. \\
\texttt{streamlit} & Rapid dashboard development & Climate-risk property
dashboards and analytic prototypes. \\
\texttt{dash} & Advanced dashboard development & More configurable
dashboard applications where needed. \\
\texttt{jupyter} & Interactive notebooks & Teaching, experimentation,
EDA and model development. \\
Quarto workflow & Reproducible reporting & Integrating code, analysis,
tables, figures and interpretation into HTML/PDF outputs. \\
Git/GitHub workflow & Reproducibility and professional portfolio &
Version control, collaboration, documentation and project
publication. \\
\end{longtable}
}

\section{Full 12-Week Curriculum Map}\label{full-12-week-curriculum-map}

\clearpage
\begin{landscape}
\begingroup
\scriptsize
\setlength{\tabcolsep}{2pt}
\renewcommand{\arraystretch}{0.92}

{\def\LTcaptype{none} % do not increment counter
\begin{longtable}[]{@{}
  >{\raggedleft\arraybackslash}p{(\linewidth - 20\tabcolsep) * \real{0.1176}}
  >{\raggedright\arraybackslash}p{(\linewidth - 20\tabcolsep) * \real{0.0882}}
  >{\raggedright\arraybackslash}p{(\linewidth - 20\tabcolsep) * \real{0.0882}}
  >{\raggedright\arraybackslash}p{(\linewidth - 20\tabcolsep) * \real{0.0882}}
  >{\raggedright\arraybackslash}p{(\linewidth - 20\tabcolsep) * \real{0.0882}}
  >{\raggedright\arraybackslash}p{(\linewidth - 20\tabcolsep) * \real{0.0882}}
  >{\raggedright\arraybackslash}p{(\linewidth - 20\tabcolsep) * \real{0.0882}}
  >{\raggedright\arraybackslash}p{(\linewidth - 20\tabcolsep) * \real{0.0882}}
  >{\raggedright\arraybackslash}p{(\linewidth - 20\tabcolsep) * \real{0.0882}}
  >{\raggedright\arraybackslash}p{(\linewidth - 20\tabcolsep) * \real{0.0882}}
  >{\raggedright\arraybackslash}p{(\linewidth - 20\tabcolsep) * \real{0.0882}}@{}}
\toprule\noalign{}
\begin{minipage}[b]{\linewidth}\raggedleft
Week
\end{minipage} & \begin{minipage}[b]{\linewidth}\raggedright
Theme
\end{minipage} & \begin{minipage}[b]{\linewidth}\raggedright
Main Topic
\end{minipage} & \begin{minipage}[b]{\linewidth}\raggedright
Learning Outcomes
\end{minipage} & \begin{minipage}[b]{\linewidth}\raggedright
Python Concepts
\end{minipage} & \begin{minipage}[b]{\linewidth}\raggedright
Domain Focus
\end{minipage} & \begin{minipage}[b]{\linewidth}\raggedright
Recommended Libraries
\end{minipage} & \begin{minipage}[b]{\linewidth}\raggedright
Practical Exercise
\end{minipage} & \begin{minipage}[b]{\linewidth}\raggedright
Dataset or Data Type
\end{minipage} & \begin{minipage}[b]{\linewidth}\raggedright
Independent Task
\end{minipage} & \begin{minipage}[b]{\linewidth}\raggedright
Expected Output / Portfolio Artifact
\end{minipage} \\
\midrule\noalign{}
\endhead
\bottomrule\noalign{}
\endlastfoot
1 & Orientation and setup & Python, notebooks, Quarto and real estate
climate-risk framing & Set up environment; understand workflow; connect
Python to domain problems & Scripts, notebooks, cells, comments, basic
expressions & Why climate risk matters for pricing, valuation and
investment & Python, JupyterLab, VS Code, Quarto, Git & Create a project
folder and first notebook analysing a small property-risk table & Small
synthetic property and risk dataset & Write a one-page analytic problem
statement & Project repository skeleton and Quarto orientation note \\
2 & Python foundations & Variables, data types, functions, loops,
conditionals & Write basic Python logic for property and climate-risk
variables & Strings, numbers, lists, dictionaries, functions, loops,
conditionals & Encoding property attributes, risk flags and investment
categories & Python standard library & Create functions to classify
property risk and investment status & Example property records & Build a
reusable property-risk classifier & Notebook with documented Python
foundations \\
3 & Data acquisition and cleaning & Excel, CSV, APIs, missing values,
joins and data quality & Import, inspect, clean and validate real estate
datasets & Reading files, data frames, missing values, joins, validation
checks & Property transactions, valuation records, EPC indicators and
flood-risk fields & pandas, numpy, requests & Clean a transaction
dataset and join risk indicators & Excel/CSV/API example datasets &
Create a data-quality report & Cleaned dataset and data-quality
notebook \\
4 & Data management with pandas and NumPy & Data transformation for
property, climate and financial data & Transform, merge, aggregate and
reshape datasets & Grouping, merging, filtering, pivoting, vectorised
operations & Local authority summaries, exposure rates and value bands &
pandas, numpy & Build local authority exposure and valuation summaries &
Property, climate-risk and local authority data & Produce a reusable
transformation script & Processed data table and transformation
notebook \\
5 & EDA and visualisation & Professional charts and exploratory analysis
& Produce charts and tables for analytical interpretation & Plotting,
distributions, correlations, grouped summaries & Flood-risk exposure,
price distributions and market segmentation & matplotlib, seaborn,
plotly & Visualise price distributions by flood-risk category & Property
transactions and flood-risk flags & Prepare an EDA brief for a valuation
audience & EDA notebook and visual summary \\
6 & Regression and econometrics & Real estate pricing and flood-risk
capitalisation & Estimate and interpret regression models & OLS, dummy
variables, interaction terms, diagnostics & Price discounts, flood-risk
capitalisation and valuation adjustment & statsmodels, scipy, pandas,
seaborn & Estimate a hedonic price model with flood-risk indicators &
Transaction, property attributes and flood-risk data & Write a
regression interpretation memo & Quarto regression report \\
7 & Machine learning & Prediction, classification and risk scoring &
Build supervised models and evaluate performance & Train/test split,
preprocessing, pipelines, metrics & Exposure scoring, vulnerability
classification, asset resilience & scikit-learn, pandas, numpy, plotly &
Predict high-risk assets or price discount categories & Property,
climate-risk and vulnerability indicators & Compare interpretable and
predictive models & Machine-learning notebook and model card \\
8 & Time-series analytics & REITs, housing prices, rates, inflation and
shocks & Analyse market dynamics and shock relationships & Date parsing,
time indexes, returns, rolling windows, volatility & REIT volatility,
infrastructure shocks, interest rates and climate events & pandas,
statsmodels, yfinance, matplotlib, plotly & Model REIT returns around
climate or infrastructure shock dates & REIT prices, rates, inflation,
event indicators & Create a time-series risk note & Time-series analysis
notebook \\
9 & Geospatial analytics & Mapping flood risk, exposure and local
vulnerability & Conduct spatial joins and produce maps & Geometries,
coordinate systems, spatial joins, buffers & Property location, flood
zones, local authorities, exposure & geopandas, shapely, folium,
rasterio/xarray where appropriate & Map properties against flood-risk
zones & Property coordinates, boundaries, flood layers & Produce an
exposure map for one case area & Geospatial notebook and interactive
map \\
10 & NLP and transition risk & ESG, policy, planning and market
sentiment text analytics & Extract transition-risk signals from
documents & Text cleaning, tokenisation, TF-IDF, classification & ESG
regulation, stranded assets, planning policy, disclosure & nltk, spaCy,
scikit-learn, pandas & Analyse ESG or planning documents for
transition-risk keywords & ESG reports, policy texts, planning documents
& Build a transition-risk signal dictionary & NLP notebook and
document-risk summary \\
11 & Dashboards & Climate-risk communication through Streamlit or Dash &
Build a prototype decision-support dashboard & App structure, widgets,
filters, charts, maps & Investor, valuer, insurer, policymaker and
asset-manager use cases & streamlit or dash, plotly, pandas, folium &
Build a property climate-risk dashboard prototype & Processed property,
risk and map data & Prepare dashboard documentation & Working dashboard
prototype \\
12 & Capstone integration & Reproducible workflow, GitHub, Quarto report
and presentation & Integrate data, models, maps, dashboard and reporting
& Project packaging, documentation, version control & Applied research,
consultancy and product pathway & Git, GitHub, Quarto, selected project
libraries & Present capstone workflow and findings & Learner-selected
capstone dataset & Complete final project and presentation & GitHub
repository, Quarto report and presentation \\
\end{longtable}
}

\endgroup
\end{landscape}
\clearpage

\section{Weekly Teaching Plan}\label{weekly-teaching-plan}

\subsection{Week 1: Python Orientation, Environment Setup and Domain
Framing}\label{week-1-python-orientation-environment-setup-and-domain-framing}

\textbf{Session purpose:} Establish the professional analytics
environment and frame Python as a tool for real estate climate-risk
evidence rather than generic coding.

\textbf{Key concepts:} Python workflow, Jupyter notebooks, VS Code,
Quarto, Git, project folders, reproducible reporting, climate-risk
analytics pipeline.

\textbf{Live teaching activities:}

\begin{itemize}
\tightlist
\item
  Install or verify Python, JupyterLab, VS Code, Quarto and Git.
\item
  Create a programme repository.
\item
  Explain the difference between notebooks, scripts and reports.
\item
  Demonstrate how a property climate-risk problem becomes a data
  workflow.
\end{itemize}

\textbf{Demonstration task:} Create a small table of property records
with columns such as \texttt{property\_id}, \texttt{local\_authority},
\texttt{sale\_price}, \texttt{flood\_risk\_band}, \texttt{epc\_rating}
and \texttt{investment\_status}.

\textbf{Guided lab:} Create a Jupyter notebook and a simple Quarto
document summarising a small sample of property-risk data.

\textbf{Independent practice:} Write a one-page problem statement on one
chosen use case: flood-risk capitalisation, physical risk and valuation,
transition risk, REIT shocks or climate-risk dashboard development.

\textbf{Suggested readings/resources:} Python official documentation;
JupyterLab documentation; Quarto documentation; GitHub Docs;
introductory real estate climate-risk reports selected and verified by
the instructor.

\textbf{Applied use case:} How climate-risk indicators enter property
valuation and investment decision-making.

\textbf{Deliverable:} Repository skeleton, first notebook and first
Quarto note.

\subsection{Week 2: Python Foundations for Property and Climate-Risk
Logic}\label{week-2-python-foundations-for-property-and-climate-risk-logic}

\textbf{Session purpose:} Build core programming logic using familiar
real estate and climate-risk variables.

\textbf{Key concepts:} Variables, data types, strings, numbers, lists,
dictionaries, functions, loops, conditionals and basic error awareness.

\textbf{Live teaching activities:}

\begin{itemize}
\tightlist
\item
  Define property attributes as Python variables.
\item
  Use lists and dictionaries to store property-level information.
\item
  Write simple functions for risk classification.
\item
  Use conditionals to assign risk categories.
\end{itemize}

\textbf{Demonstration task:} Write a function that classifies a property
as low, medium or high risk based on flood-risk band, EPC rating and
location vulnerability.

\textbf{Guided lab:} Create property records and loop through them to
assign investment-risk labels.

\textbf{Independent practice:} Extend the classifier to include
mortgageability, insurance affordability or infrastructure vulnerability
indicators.

\textbf{Suggested readings/resources:} Python official tutorial;
selected beginner Python chapters or open resources; instructor-prepared
property-risk coding examples.

\textbf{Applied use case:} Translating valuation and investment rules
into transparent computational logic.

\textbf{Deliverable:} Python foundations notebook with documented
functions and classification logic.

\subsection{Week 3: Data Import, Cleaning and Quality
Assurance}\label{week-3-data-import-cleaning-and-quality-assurance}

\textbf{Session purpose:} Teach the learner to bring real estate and
climate-risk data into Python and prepare it for analysis.

\textbf{Key concepts:} CSV and Excel import, API basics, column naming,
missing values, duplicates, data types, joins, validation checks and
data dictionaries.

\textbf{Live teaching activities:}

\begin{itemize}
\tightlist
\item
  Import Excel and CSV files.
\item
  Diagnose missing values and inconsistent categories.
\item
  Convert date, numeric and categorical fields.
\item
  Merge transaction data with risk indicators.
\item
  Create a data-quality checklist.
\end{itemize}

\textbf{Demonstration task:} Clean a property transaction dataset and
merge it with a flood-risk indicator table.

\textbf{Guided lab:} Produce a data-quality report showing missing
values, duplicate records, invalid price values and unmatched local
authorities.

\textbf{Independent practice:} Prepare a cleaned version of a
property-climate dataset and document all cleaning decisions.

\textbf{Suggested readings/resources:} pandas documentation on
input/output, missing data and merging; documentation for relevant
open-data portals, to be verified locally.

\textbf{Applied use case:} Preparing transaction and flood-risk data for
pricing analysis.

\textbf{Deliverable:} Cleaned dataset and data-quality notebook.

\subsection{Week 4: Data Management with pandas and
NumPy}\label{week-4-data-management-with-pandas-and-numpy}

\textbf{Session purpose:} Build fluency in transforming and summarising
large property, climate-risk and financial datasets.

\textbf{Key concepts:} Filtering, selecting, grouping, aggregation,
pivot tables, reshaping, vectorised operations and feature construction.

\textbf{Live teaching activities:}

\begin{itemize}
\tightlist
\item
  Create derived variables such as price per square metre, risk score,
  EPC category and exposure dummy.
\item
  Aggregate exposure by local authority.
\item
  Compare average prices across risk bands.
\item
  Build reusable transformation steps.
\end{itemize}

\textbf{Demonstration task:} Create a local authority summary table
showing number of properties, mean price, flood exposure rate, average
EPC score and risk category.

\textbf{Guided lab:} Transform raw property and risk datasets into an
analysis-ready table.

\textbf{Independent practice:} Write a short script that takes raw input
and produces a processed dataset.

\textbf{Suggested readings/resources:} pandas user guide; NumPy
quickstart; instructor-prepared data transformation checklist.

\textbf{Applied use case:} Turning fragmented property and risk records
into an investment-ready analytical dataset.

\textbf{Deliverable:} Processed dataset and transformation notebook or
script.

\subsection{Week 5: Exploratory Data Analysis and Professional
Visualisation}\label{week-5-exploratory-data-analysis-and-professional-visualisation}

\textbf{Session purpose:} Develop the ability to explore and communicate
data patterns before formal modelling.

\textbf{Key concepts:} Summary statistics, distributions, outliers,
grouped comparisons, correlations, static and interactive charts.

\textbf{Live teaching activities:}

\begin{itemize}
\tightlist
\item
  Produce price distribution plots.
\item
  Compare exposed and non-exposed properties.
\item
  Create correlation heatmaps.
\item
  Use interactive charts for stakeholder communication.
\item
  Discuss visual ethics and avoiding misleading graphics.
\end{itemize}

\textbf{Demonstration task:} Visualise sale-price distributions across
flood-risk bands and EPC categories.

\textbf{Guided lab:} Create a professional EDA notebook for a valuation
or investment audience.

\textbf{Independent practice:} Prepare a two-page visual brief
summarising one dataset's climate-risk patterns.

\textbf{Suggested readings/resources:} Matplotlib documentation; seaborn
documentation; Plotly documentation; selected examples of professional
real estate market reports.

\textbf{Applied use case:} Understanding whether flood-risk exposure
appears associated with observed price differences before formal
modelling.

\textbf{Deliverable:} EDA notebook and visual summary.

\subsection{Week 6: Regression and Econometric Modelling for Real Estate
Pricing}\label{week-6-regression-and-econometric-modelling-for-real-estate-pricing}

\textbf{Session purpose:} Apply statistical and econometric modelling to
real estate pricing and climate-risk capitalisation.

\textbf{Key concepts:} OLS regression, dependent and independent
variables, dummy variables, controls, interactions, model assumptions,
diagnostics and interpretation.

\textbf{Live teaching activities:}

\begin{itemize}
\tightlist
\item
  Specify a hedonic pricing model.
\item
  Estimate regression using \texttt{statsmodels}.
\item
  Interpret coefficients for flood-risk bands and property attributes.
\item
  Discuss omitted variable bias, spatial confounding and model
  limitations.
\item
  Produce regression tables suitable for reports.
\end{itemize}

\textbf{Demonstration task:} Estimate a model of sale price or log sale
price using property attributes, flood-risk indicators, EPC rating and
local authority controls.

\textbf{Guided lab:} Build and interpret a flood-risk capitalisation
regression model.

\textbf{Independent practice:} Write a regression interpretation memo
for a professional valuation or investment audience.

\textbf{Suggested readings/resources:} statsmodels documentation;
instructor-selected readings on hedonic pricing and climate-risk
capitalisation, to be verified for the specific context.

\textbf{Applied use case:} Estimating whether flood-risk exposure is
associated with a price discount after controlling for property
characteristics.

\textbf{Deliverable:} Quarto regression report.

\subsection{Week 7: Machine Learning for Prediction, Classification and
Risk
Scoring}\label{week-7-machine-learning-for-prediction-classification-and-risk-scoring}

\textbf{Session purpose:} Introduce machine learning as a complement to
interpretable statistical modelling.

\textbf{Key concepts:} Supervised learning, train/test split,
preprocessing, pipelines, cross-validation, evaluation metrics, feature
importance, model cards and responsible interpretation.

\textbf{Live teaching activities:}

\begin{itemize}
\tightlist
\item
  Define prediction and classification tasks.
\item
  Build baseline models.
\item
  Compare linear models, tree-based models and simple ensemble
  approaches.
\item
  Discuss overfitting and model interpretability.
\item
  Create a model card documenting purpose, data, limitations and
  intended use.
\end{itemize}

\textbf{Demonstration task:} Predict whether a property belongs to a
high-risk investment category based on location, flood-risk, EPC, price
and vulnerability indicators.

\textbf{Guided lab:} Build and evaluate a machine-learning model for
exposure scoring or price-discount classification.

\textbf{Independent practice:} Compare a transparent model with a more
predictive model and explain the trade-offs.

\textbf{Suggested readings/resources:} scikit-learn documentation; model
evaluation guides; responsible AI and model documentation resources
selected by the instructor.

\textbf{Applied use case:} Classifying assets into climate-risk
investment bands for preliminary screening.

\textbf{Deliverable:} Machine-learning notebook and model card.

\subsection{Week 8: Time-Series Analysis for REITs, Housing Prices and
Market
Shocks}\label{week-8-time-series-analysis-for-reits-housing-prices-and-market-shocks}

\textbf{Session purpose:} Develop skills for analysing temporal patterns
in financial and real estate market data.

\textbf{Key concepts:} Date-time parsing, time indexes, returns,
volatility, rolling windows, event indicators, lagged variables and
basic time-series models.

\textbf{Live teaching activities:}

\begin{itemize}
\tightlist
\item
  Import and clean financial time-series data.
\item
  Calculate REIT returns and rolling volatility.
\item
  Align shock dates with market data.
\item
  Compare periods before and after climate, energy, interest-rate or
  infrastructure events.
\item
  Discuss correlation, causality and event-study limitations.
\end{itemize}

\textbf{Demonstration task:} Analyse REIT return volatility around an
infrastructure shock, policy uncertainty period or climate-related
event.

\textbf{Guided lab:} Create time-series plots and basic regression/event
indicators for REIT or housing price data.

\textbf{Independent practice:} Prepare a time-series risk note linking
market behaviour to one selected shock type.

\textbf{Suggested readings/resources:} pandas time-series documentation;
statsmodels time-series documentation; yfinance documentation or
alternative finance-data source documentation, subject to verification
and data-use terms.

\textbf{Applied use case:} Understanding how climate events,
infrastructure disruption or policy uncertainty may affect REIT returns
and volatility.

\textbf{Deliverable:} Time-series analysis notebook and risk note.

\subsection{Week 9: Geospatial Analytics for Flood Exposure and Local
Vulnerability}\label{week-9-geospatial-analytics-for-flood-exposure-and-local-vulnerability}

\textbf{Session purpose:} Enable the learner to integrate property
location, flood-risk layers and local authority indicators.

\textbf{Key concepts:} Coordinate reference systems, points, polygons,
shapefiles, GeoJSON, spatial joins, buffers, choropleth maps,
interactive mapping and raster awareness.

\textbf{Live teaching activities:}

\begin{itemize}
\tightlist
\item
  Load boundary and property-location data.
\item
  Convert latitude/longitude data into geometries.
\item
  Join property points to local authority polygons.
\item
  Overlay properties with flood-risk zones.
\item
  Build interactive exposure maps.
\end{itemize}

\textbf{Demonstration task:} Map property exposure to flood-risk zones
and summarise exposure by local authority.

\textbf{Guided lab:} Conduct spatial joins and produce an interactive
flood-risk map.

\textbf{Independent practice:} Create a short geospatial exposure report
for one selected area.

\textbf{Suggested readings/resources:} GeoPandas documentation; Shapely
documentation; Folium documentation; relevant national or local
geospatial data portals, to be verified.

\textbf{Applied use case:} Identifying property assets exposed to flood
risk and linking exposure to valuation and vulnerability indicators.

\textbf{Deliverable:} Geospatial notebook and interactive flood-risk
map.

\subsection{Week 10: NLP for ESG, Planning, Policy and Transition-Risk
Signals}\label{week-10-nlp-for-esg-planning-policy-and-transition-risk-signals}

\textbf{Session purpose:} Introduce text analytics for transition-risk
evidence in real estate investment and policy contexts.

\textbf{Key concepts:} Text extraction, cleaning, tokenisation, stopword
removal, keyword dictionaries, TF-IDF, topic exploration, simple
classification and limitations of automated text analysis.

\textbf{Live teaching activities:}

\begin{itemize}
\tightlist
\item
  Load a small corpus of ESG, planning or policy documents.
\item
  Clean and tokenize text.
\item
  Build a transition-risk keyword dictionary.
\item
  Extract document-level signals.
\item
  Discuss interpretive validity and the need for human review.
\end{itemize}

\textbf{Demonstration task:} Analyse ESG disclosures or policy documents
for signals related to energy efficiency, regulation, carbon
performance, stranded assets and retrofit risk.

\textbf{Guided lab:} Build a document-risk summary table for a small
corpus.

\textbf{Independent practice:} Compare transition-risk signals across
document types.

\textbf{Suggested readings/resources:} nltk documentation; spaCy
documentation; scikit-learn text feature extraction documentation;
selected ESG or planning documents, to be verified for access and use
rights.

\textbf{Applied use case:} Identifying transition-risk signals relevant
to real estate investors from policy and disclosure texts.

\textbf{Deliverable:} NLP notebook and transition-risk document summary.

\subsection{Week 11: Dashboard Development for Real Estate Climate-Risk
Communication}\label{week-11-dashboard-development-for-real-estate-climate-risk-communication}

\textbf{Session purpose:} Build a prototype decision-support dashboard
for communicating climate-risk analytics.

\textbf{Key concepts:} Dashboard structure, user stories, filters,
inputs, charts, maps, tables, layout, interactivity and deployment
basics.

\textbf{Live teaching activities:}

\begin{itemize}
\tightlist
\item
  Define dashboard users: investor, valuer, insurer, policymaker, local
  authority and asset manager.
\item
  Build a simple Streamlit or Dash app.
\item
  Add filters for local authority, flood-risk band, EPC rating and
  property value.
\item
  Add visual summaries and maps.
\item
  Discuss dashboard limitations and data governance.
\end{itemize}

\textbf{Demonstration task:} Build a prototype dashboard showing
property value, exposure, EPC rating, vulnerability indicators and
investment-risk score.

\textbf{Guided lab:} Convert a processed dataset into an interactive
dashboard.

\textbf{Independent practice:} Write dashboard documentation explaining
data inputs, assumptions, limitations and intended users.

\textbf{Suggested readings/resources:} Streamlit documentation; Dash
documentation; Plotly documentation; dashboard design examples selected
by the instructor.

\textbf{Applied use case:} Communicating asset-level climate-risk
evidence to investors, valuers, insurers, policymakers and asset
managers.

\textbf{Deliverable:} Working dashboard prototype and documentation.

\subsection{Week 12: Capstone Integration, Reproducibility and
Professional
Presentation}\label{week-12-capstone-integration-reproducibility-and-professional-presentation}

\textbf{Session purpose:} Integrate the programme into a final applied
project with documented workflow, professional outputs and future
development pathway.

\textbf{Key concepts:} Capstone synthesis, GitHub repository, Quarto
technical report, reproducible workflow, presentation, model limitations
and stakeholder communication.

\textbf{Live teaching activities:}

\begin{itemize}
\tightlist
\item
  Review capstone requirements.
\item
  Conduct repository and documentation checks.
\item
  Refine analysis, visuals, maps and dashboards.
\item
  Prepare a professional presentation.
\item
  Discuss publication, consultancy and product-development extensions.
\end{itemize}

\textbf{Demonstration task:} Walk through a complete end-to-end project
structure from raw data to dashboard and Quarto report.

\textbf{Guided lab:} Capstone troubleshooting and presentation
rehearsal.

\textbf{Independent practice:} Finalise the capstone report,
dashboard/prototype and presentation.

\textbf{Suggested readings/resources:} Quarto documentation; GitHub
Docs; project-specific technical documentation and verified domain
resources.

\textbf{Applied use case:} Producing a research, consultancy or
product-ready analytical workflow.

\textbf{Deliverable:} Final capstone repository, Quarto report and
presentation.

\section{Applied Dataset Strategy}\label{applied-dataset-strategy}

The programme should use datasets that reflect the learner's intended
research and professional work. Some datasets may be public, some
institutional, some commercial and some locally sourced. Where specific
access is uncertain, dataset references should be treated as examples to
be verified, adapted to local availability and checked against data-use
rights.

{\def\LTcaptype{none} % do not increment counter
\begin{longtable}[]{@{}
  >{\raggedright\arraybackslash}p{(\linewidth - 6\tabcolsep) * \real{0.2500}}
  >{\raggedright\arraybackslash}p{(\linewidth - 6\tabcolsep) * \real{0.2500}}
  >{\raggedright\arraybackslash}p{(\linewidth - 6\tabcolsep) * \real{0.2500}}
  >{\raggedright\arraybackslash}p{(\linewidth - 6\tabcolsep) * \real{0.2500}}@{}}
\toprule\noalign{}
\begin{minipage}[b]{\linewidth}\raggedright
Dataset Type
\end{minipage} & \begin{minipage}[b]{\linewidth}\raggedright
Possible Variables
\end{minipage} & \begin{minipage}[b]{\linewidth}\raggedright
Use in Curriculum
\end{minipage} & \begin{minipage}[b]{\linewidth}\raggedright
Access Note
\end{minipage} \\
\midrule\noalign{}
\endhead
\bottomrule\noalign{}
\endlastfoot
Property transaction data & Sale price, date, property type, tenure,
location, floor area, age, transaction identifier & Pricing models, EDA,
flood-risk capitalisation & Public, commercial or institutional source
to be verified. \\
Housing price indices & Area, time period, index value, growth rate &
Time-series analysis and market context & National statistical or
property-market source to be verified. \\
Valuation data & Assessed value, valuation date, property attributes,
adjustment factors & Valuation modelling and risk adjustment & Often
commercial or confidential; use anonymised or synthetic examples if
needed. \\
Flood-risk indicators & Flood zone, risk band, return period, exposure
score, defence status & Flood-risk capitalisation, exposure mapping,
dashboards & Environmental agency or local authority data to be
verified. \\
Climate exposure data & Heat, flood, storm, sea-level, precipitation or
hazard indicators & Physical climate-risk analytics & Source depends on
jurisdiction and scale; verify suitability. \\
EPC or energy-efficiency data & EPC rating, energy consumption,
emissions estimate, certificate date & Transition-risk and
stranded-asset analysis & Public or regulated data source depending on
country. \\
REIT price and returns data & Price, adjusted close, returns, volume,
market index & REIT volatility and shock analysis & Financial data
provider or open finance API, subject to terms. \\
Interest rates & Policy rate, mortgage rate, bond yield & Time-series
and investment-risk modelling & Central bank or financial data source to
be verified. \\
Inflation data & CPI, inflation rate, construction cost index & Market
stress and valuation context & National statistical source to be
verified. \\
Infrastructure shock data & Power outages, transport disruptions,
utility failures, event dates & Event analysis and infrastructure
resilience & Local authority, utility or news-derived source to be
verified. \\
ESG reports & Company disclosures, sustainability reports, transition
plans & NLP and transition-risk signal extraction & Company websites or
disclosure databases, subject to use rights. \\
Policy documents & Climate policy, planning regulation, building
standards, retrofit rules & NLP and transition-risk analysis &
Government or regulator sources to be verified. \\
Planning documents & Local plans, development frameworks, flood-planning
guidance & Policy and market-sentiment text analytics & Local authority
sources to be verified. \\
Local authority vulnerability indicators & Income, deprivation,
infrastructure, housing stock, demographic indicators & Vulnerability
scoring and spatial analysis & Statistical or local government source to
be verified. \\
Insurance affordability indicators & Premiums, coverage availability,
claims exposure, insurer withdrawal & Mortgageability and liquidity
analysis & Often restricted; use proxy indicators if direct data
unavailable. \\
Mortgageability indicators & Lending restrictions, loan-to-value
constraints, risk flags & Investment and liquidity analysis & Commercial
or institutional source to be verified. \\
\end{longtable}
}

\section{Practical Exercises}\label{practical-exercises}

{\def\LTcaptype{none} % do not increment counter
\begin{longtable}[]{@{}
  >{\raggedleft\arraybackslash}p{(\linewidth - 6\tabcolsep) * \real{0.3077}}
  >{\raggedright\arraybackslash}p{(\linewidth - 6\tabcolsep) * \real{0.2308}}
  >{\raggedright\arraybackslash}p{(\linewidth - 6\tabcolsep) * \real{0.2308}}
  >{\raggedright\arraybackslash}p{(\linewidth - 6\tabcolsep) * \real{0.2308}}@{}}
\toprule\noalign{}
\begin{minipage}[b]{\linewidth}\raggedleft
Week
\end{minipage} & \begin{minipage}[b]{\linewidth}\raggedright
Exercise Title
\end{minipage} & \begin{minipage}[b]{\linewidth}\raggedright
Task
\end{minipage} & \begin{minipage}[b]{\linewidth}\raggedright
Expected Output
\end{minipage} \\
\midrule\noalign{}
\endhead
\bottomrule\noalign{}
\endlastfoot
1 & Build the Real Estate Climate-Risk Project Shell & Create a Python
project folder, notebook, Quarto report and README for one selected use
case. & Project folder and first Quarto note. \\
2 & Encode Property-Risk Logic in Python & Write functions that classify
property climate-risk status using flood risk, EPC rating and
vulnerability indicators. & Python foundations notebook. \\
3 & Clean a Property Transaction Dataset & Import Excel/CSV records,
correct data types, handle missing values and merge with a risk table. &
Cleaned dataset and data-quality report. \\
4 & Create Local Authority Risk Summaries & Aggregate property values,
exposure rates and EPC scores by local authority. & Processed summary
table. \\
5 & Visualise Flood-Risk Capitalisation Patterns & Produce charts
comparing price distributions and exposure categories. & EDA visual
report. \\
6 & Estimate a Hedonic Pricing Model & Run a regression model linking
sale price to property attributes and flood-risk indicators. &
Regression report in Quarto. \\
7 & Build an Investment-Risk Classifier & Train a model to classify
properties into investment-risk categories. & Machine-learning notebook
and model card. \\
8 & Analyse REIT Volatility Around Shocks & Calculate returns and
rolling volatility around climate, infrastructure or policy events. &
Time-series risk note. \\
9 & Map Property Flood Exposure & Conduct a spatial join between
property points and flood-risk polygons. & Interactive flood-risk
map. \\
10 & Extract Transition-Risk Signals from Text & Analyse ESG, planning
or policy documents for transition-risk indicators. & NLP summary table
and interpretation. \\
11 & Build a Climate-Risk Dashboard Prototype & Create a Streamlit or
Dash dashboard with filters, charts and maps. & Working dashboard
prototype. \\
12 & Present an End-to-End Analytical System & Integrate data, model,
map, dashboard and report into a capstone presentation. & Final
repository, report and presentation. \\
\end{longtable}
}

\section{Assessment Design}\label{assessment-design}

The assessment structure should be practical, cumulative and
portfolio-oriented.

{\def\LTcaptype{none} % do not increment counter
\begin{longtable}[]{@{}
  >{\raggedright\arraybackslash}p{(\linewidth - 6\tabcolsep) * \real{0.2308}}
  >{\raggedleft\arraybackslash}p{(\linewidth - 6\tabcolsep) * \real{0.3077}}
  >{\raggedright\arraybackslash}p{(\linewidth - 6\tabcolsep) * \real{0.2308}}
  >{\raggedright\arraybackslash}p{(\linewidth - 6\tabcolsep) * \real{0.2308}}@{}}
\toprule\noalign{}
\begin{minipage}[b]{\linewidth}\raggedright
Assessment Component
\end{minipage} & \begin{minipage}[b]{\linewidth}\raggedleft
Weight
\end{minipage} & \begin{minipage}[b]{\linewidth}\raggedright
Description
\end{minipage} & \begin{minipage}[b]{\linewidth}\raggedright
Evidence
\end{minipage} \\
\midrule\noalign{}
\endhead
\bottomrule\noalign{}
\endlastfoot
Weekly coding exercises & 20\% & Short weekly exercises demonstrating
incremental Python skills in applied real estate and climate-risk
contexts. & Notebooks, scripts and outputs. \\
Short applied analytics tasks & 15\% & Brief analytical memos or outputs
interpreting data patterns for valuation, investment, policy or risk
audiences. & Visual briefs, risk notes and interpretation memos. \\
Mid-programme mini-project & 15\% & A small integrated project combining
data cleaning, EDA and regression or preliminary modelling. & Quarto
mini-report and cleaned dataset. \\
GitHub / reproducibility portfolio & 15\% & Evidence of organised
project folders, README files, version control, documentation and
reproducible workflow. & GitHub repository or local repository
export. \\
Final capstone project & 25\% & End-to-end applied project selected from
the capstone options. & Data workflow, model, map/dashboard and Quarto
report. \\
Professional presentation or report defence & 10\% & Concise
presentation of the capstone problem, method, findings, limitations and
stakeholder relevance. & Slide deck, oral presentation or recorded
walkthrough. \\
\textbf{Total} & \textbf{100\%} & & \\
\end{longtable}
}

\section{Assessment Rubric}\label{assessment-rubric}

{\def\LTcaptype{none} % do not increment counter
\begin{longtable}[]{@{}
  >{\raggedright\arraybackslash}p{(\linewidth - 8\tabcolsep) * \real{0.2000}}
  >{\raggedright\arraybackslash}p{(\linewidth - 8\tabcolsep) * \real{0.2000}}
  >{\raggedright\arraybackslash}p{(\linewidth - 8\tabcolsep) * \real{0.2000}}
  >{\raggedright\arraybackslash}p{(\linewidth - 8\tabcolsep) * \real{0.2000}}
  >{\raggedright\arraybackslash}p{(\linewidth - 8\tabcolsep) * \real{0.2000}}@{}}
\toprule\noalign{}
\begin{minipage}[b]{\linewidth}\raggedright
Criterion
\end{minipage} & \begin{minipage}[b]{\linewidth}\raggedright
Excellent
\end{minipage} & \begin{minipage}[b]{\linewidth}\raggedright
Good
\end{minipage} & \begin{minipage}[b]{\linewidth}\raggedright
Developing
\end{minipage} & \begin{minipage}[b]{\linewidth}\raggedright
Needs Improvement
\end{minipage} \\
\midrule\noalign{}
\endhead
\bottomrule\noalign{}
\endlastfoot
Technical accuracy & Code runs cleanly, is well structured and produces
correct outputs. & Code mostly runs with minor inefficiencies or small
issues. & Code runs partially but requires substantial correction. &
Code does not run or outputs are unreliable. \\
Data cleaning and management & Cleaning decisions are systematic,
documented and appropriate to the dataset. & Most cleaning steps are
appropriate and understandable. & Some cleaning steps are attempted but
inconsistently applied. & Data cleaning is weak, undocumented or
incorrect. \\
Analytical reasoning & Analysis is clearly linked to a real estate
climate-risk problem and uses appropriate logic. & Analysis is relevant
but could be more tightly justified. & Analysis is partially relevant
but underdeveloped. & Analysis lacks clear connection to the problem. \\
Model appropriateness & Model choice is justified, limitations are
stated and outputs are interpreted responsibly. & Model choice is
generally suitable with some limitations discussed. & Model is attempted
but justification or interpretation is weak. & Model is inappropriate,
unexplained or misinterpreted. \\
Domain interpretation & Findings are interpreted in relation to
valuation, investment, flood risk, physical risk, transition risk or
policy relevance. & Interpretation is relevant but not fully developed.
& Interpretation is descriptive rather than analytical. & Interpretation
is absent or disconnected from the domain. \\
Visualisation quality & Charts, maps and tables are clear, professional,
labelled and decision-relevant. & Visuals are understandable with minor
design issues. & Visuals are basic or inconsistently labelled. & Visuals
are unclear, misleading or absent. \\
Reproducibility & Project has clear folder structure, documented
dependencies, Quarto report and version-control evidence. &
Reproducibility is mostly evident but documentation could improve. &
Some reproducibility practices are present. & Workflow is disorganised
and difficult to reproduce. \\
Documentation & README, comments, data notes and method explanations are
clear and professional. & Documentation is present but incomplete in
places. & Documentation is minimal or inconsistent. & Documentation is
missing or unclear. \\
Policy / investment relevance & Output clearly addresses stakeholder
decisions and professional implications. & Relevance is visible but
could be more explicit. & Relevance is implied but weakly articulated. &
Stakeholder relevance is absent. \\
Communication & Report or presentation is polished, coherent and
suitable for professional or academic audiences. & Communication is
clear with minor structural weaknesses. & Communication is
understandable but uneven. & Communication is unclear or incomplete. \\
\end{longtable}
}

\section{Capstone Project Options}\label{capstone-project-options}

\subsection{Capstone A: Flood Risk and House Price Analytics
Tool}\label{capstone-a-flood-risk-and-house-price-analytics-tool}

\subsubsection{Project rationale}\label{project-rationale}

Flood risk can affect housing prices, market liquidity, insurance
affordability, mortgageability and valuation confidence. This project
develops an analytical workflow for examining whether and how flood-risk
indicators are associated with house prices.

\subsubsection{Research / professional
problem}\label{research-professional-problem}

Can Python be used to estimate and communicate the relationship between
flood-risk exposure and house prices in a way that is useful for
researchers, valuers, investors, insurers, policymakers and local
authorities?

\subsubsection{Guiding questions}\label{guiding-questions}

\begin{itemize}
\tightlist
\item
  Are properties in higher flood-risk categories associated with lower
  observed prices?
\item
  Does the relationship vary by property type, location or local
  authority?
\item
  How sensitive are results to model specification?
\item
  What are the implications for valuation, investment risk and market
  liquidity?
\end{itemize}

\subsubsection{Required data types}\label{required-data-types}

\begin{itemize}
\tightlist
\item
  Property transaction data.
\item
  Property attributes.
\item
  Flood-risk indicators.
\item
  Local authority identifiers or boundaries.
\item
  Optional: EPC rating, vulnerability indicators, insurance proxy
  variables.
\end{itemize}

\subsubsection{Python libraries}\label{python-libraries}

\texttt{pandas}, \texttt{numpy}, \texttt{matplotlib}, \texttt{seaborn},
\texttt{plotly}, \texttt{statsmodels}, \texttt{geopandas},
\texttt{shapely}, \texttt{folium}, \texttt{quarto}.

\subsubsection{Analytic workflow}\label{analytic-workflow}

\begin{enumerate}
\def\labelenumi{\arabic{enumi}.}
\tightlist
\item
  Define the research question and case area.
\item
  Import transaction and flood-risk datasets.
\item
  Clean property records and standardise variables.
\item
  Join property and risk data.
\item
  Conduct EDA by risk band and location.
\item
  Estimate hedonic regression models.
\item
  Interpret flood-risk coefficients cautiously.
\item
  Map exposed properties and local authority patterns.
\item
  Prepare a Quarto report.
\item
  Package the workflow in a documented repository.
\end{enumerate}

\subsubsection{Expected outputs}\label{expected-outputs}

\begin{itemize}
\tightlist
\item
  Cleaned analysis-ready dataset.
\item
  EDA charts.
\item
  Regression model and interpretation.
\item
  Flood-risk exposure map.
\item
  Quarto technical report.
\item
  Professional summary memo.
\end{itemize}

\subsubsection{Dashboard / report
expectations}\label{dashboard-report-expectations}

The report should include the research problem, data description,
cleaning decisions, model specification, key findings, visual evidence,
limitations and implications for valuation or investment
decision-making.

\subsubsection{Assessment criteria}\label{assessment-criteria}

\begin{itemize}
\tightlist
\item
  Quality of data cleaning.
\item
  Appropriateness of regression specification.
\item
  Carefulness of interpretation.
\item
  Clarity of charts and maps.
\item
  Reproducibility of the workflow.
\item
  Relevance to valuation and investment decision-making.
\end{itemize}

\subsubsection{Extension potential}\label{extension-potential}

The project can be extended into a journal article on flood-risk
capitalisation, a consultancy tool for valuers, a grant-funded urban
resilience study or an investor-facing flood-risk screening product.

\subsection{Capstone B: Climate-Risk Property
Dashboard}\label{capstone-b-climate-risk-property-dashboard}

\subsubsection{Project rationale}\label{project-rationale-1}

Stakeholders need accessible tools for understanding asset-level and
area-level climate-risk exposure. A dashboard can translate complex
property, climate, vulnerability and investment indicators into
decision-support evidence.

\subsubsection{Research / professional
problem}\label{research-professional-problem-1}

How can Python be used to build an interactive dashboard that
communicates flood exposure, EPC rating, property value, local authority
context, vulnerability indicators and investment-risk scoring?

\subsubsection{Guiding questions}\label{guiding-questions-1}

\begin{itemize}
\tightlist
\item
  Which assets are most exposed to climate-risk indicators?
\item
  How do flood risk, EPC rating and property value interact?
\item
  Which local authorities show higher vulnerability or investment
  concern?
\item
  How can dashboard design support investors, valuers, insurers,
  policymakers and asset managers?
\end{itemize}

\subsubsection{Required data types}\label{required-data-types-1}

\begin{itemize}
\tightlist
\item
  Property-level dataset.
\item
  Flood-risk indicators.
\item
  EPC or energy-efficiency data.
\item
  Local authority data.
\item
  Vulnerability indicators.
\item
  Optional: valuation, mortgageability or insurance proxy indicators.
\end{itemize}

\subsubsection{Python libraries}\label{python-libraries-1}

\texttt{pandas}, \texttt{numpy}, \texttt{plotly}, \texttt{geopandas},
\texttt{folium}, \texttt{streamlit} or \texttt{dash}, \texttt{shapely}.

\subsubsection{Analytic workflow}\label{analytic-workflow-1}

\begin{enumerate}
\def\labelenumi{\arabic{enumi}.}
\tightlist
\item
  Define dashboard users and decisions.
\item
  Create a cleaned dashboard dataset.
\item
  Build a risk-scoring logic.
\item
  Create charts, maps and summary tables.
\item
  Build filters by location, risk band, EPC rating and value band.
\item
  Add explanatory text and limitations.
\item
  Test usability with a simple user story.
\item
  Document the dashboard and repository.
\end{enumerate}

\subsubsection{Expected outputs}\label{expected-outputs-1}

\begin{itemize}
\tightlist
\item
  Dashboard-ready dataset.
\item
  Risk-scoring method note.
\item
  Interactive dashboard.
\item
  Map of property exposure.
\item
  README and user guide.
\end{itemize}

\subsubsection{Dashboard / report
expectations}\label{dashboard-report-expectations-1}

The dashboard should include filters, charts, summary metrics, property
or area-level risk indicators and clear methodological notes. It must
avoid presenting risk scores as definitive without explaining
assumptions.

\subsubsection{Assessment criteria}\label{assessment-criteria-1}

\begin{itemize}
\tightlist
\item
  Dashboard functionality.
\item
  Relevance of indicators.
\item
  Clarity of user interface.
\item
  Transparency of risk scoring.
\item
  Quality of documentation.
\item
  Suitability for stakeholder communication.
\end{itemize}

\subsubsection{Extension potential}\label{extension-potential-1}

This project can become a consultancy prototype, policy-facing tool,
investor dashboard, teaching demonstration or foundation for
grant-funded digital infrastructure work.

\subsection{Capstone C: REIT Climate and Infrastructure Shock
Model}\label{capstone-c-reit-climate-and-infrastructure-shock-model}

\subsubsection{Project rationale}\label{project-rationale-2}

REITs and listed real estate vehicles are exposed to climate events,
infrastructure disruption, energy shocks, policy uncertainty,
interest-rate shifts and wider financial-market volatility. Python can
support analysis of how these shocks relate to returns and volatility.

\subsubsection{Research / professional
problem}\label{research-professional-problem-2}

How do climate, infrastructure, policy or macro-financial shocks
correspond with REIT returns, volatility and investment risk?

\subsubsection{Guiding questions}\label{guiding-questions-2}

\begin{itemize}
\tightlist
\item
  How do REIT returns behave around selected shock dates?
\item
  Does volatility increase during climate, infrastructure or policy
  events?
\item
  How do interest rates and inflation interact with REIT performance?
\item
  What are the limitations of attributing financial-market movements to
  climate or infrastructure shocks?
\end{itemize}

\subsubsection{Required data types}\label{required-data-types-2}

\begin{itemize}
\tightlist
\item
  REIT price data.
\item
  Market index data.
\item
  Interest-rate data.
\item
  Inflation data.
\item
  Climate-event or infrastructure-shock dates.
\item
  Optional: energy-price or policy-uncertainty indicators.
\end{itemize}

\subsubsection{Python libraries}\label{python-libraries-2}

\texttt{pandas}, \texttt{numpy}, \texttt{matplotlib}, \texttt{plotly},
\texttt{statsmodels}, \texttt{yfinance} or alternative finance-data
package, \texttt{scipy}.

\subsubsection{Analytic workflow}\label{analytic-workflow-2}

\begin{enumerate}
\def\labelenumi{\arabic{enumi}.}
\tightlist
\item
  Select REITs or listed real estate securities.
\item
  Import price and macro-financial data.
\item
  Calculate returns and rolling volatility.
\item
  Construct shock-event indicators.
\item
  Visualise market behaviour around events.
\item
  Estimate simple models relating returns or volatility to shock
  variables.
\item
  Interpret cautiously with attention to confounding.
\item
  Produce a professional risk note.
\end{enumerate}

\subsubsection{Expected outputs}\label{expected-outputs-2}

\begin{itemize}
\tightlist
\item
  Cleaned time-series dataset.
\item
  Return and volatility charts.
\item
  Event-window summary.
\item
  Basic regression or volatility model.
\item
  Quarto risk report.
\end{itemize}

\subsubsection{Dashboard / report
expectations}\label{dashboard-report-expectations-2}

The report should include clear event definitions, data limitations,
market context, charts, model outputs and implications for investment
strategy.

\subsubsection{Assessment criteria}\label{assessment-criteria-2}

\begin{itemize}
\tightlist
\item
  Correct handling of time-series data.
\item
  Appropriate calculation of returns and volatility.
\item
  Sensible event framing.
\item
  Responsible interpretation.
\item
  Communication of uncertainty.
\item
  Reproducibility.
\end{itemize}

\subsubsection{Extension potential}\label{extension-potential-2}

The project can support research on REIT climate-risk exposure,
infrastructure resilience, financial-market vulnerability or
climate-adjusted real estate investment strategies.

\subsection{Capstone D: Transition Risk Text Analytics
Project}\label{capstone-d-transition-risk-text-analytics-project}

\subsubsection{Project rationale}\label{project-rationale-3}

Transition risk is embedded in policies, ESG disclosures, planning
documents, regulation, retrofit expectations, market sentiment and
investor communication. NLP can help identify signals that matter for
real estate assets and portfolios.

\subsubsection{Research / professional
problem}\label{research-professional-problem-3}

Can Python-based text analytics identify transition-risk signals in ESG
reports, planning documents or policy texts relevant to real estate
investors?

\subsubsection{Guiding questions}\label{guiding-questions-3}

\begin{itemize}
\tightlist
\item
  Which terms and themes indicate transition-risk exposure?
\item
  How do documents discuss energy efficiency, retrofit, carbon
  performance, regulation and stranded assets?
\item
  Do different document types emphasise different risk signals?
\item
  How should automated text analysis be combined with expert
  interpretation?
\end{itemize}

\subsubsection{Required data types}\label{required-data-types-3}

\begin{itemize}
\tightlist
\item
  ESG reports.
\item
  Planning documents.
\item
  Policy documents.
\item
  Market commentary or investor disclosures where permitted.
\item
  Transition-risk keyword dictionary.
\end{itemize}

\subsubsection{Python libraries}\label{python-libraries-3}

\texttt{pandas}, \texttt{nltk}, \texttt{spaCy}, \texttt{scikit-learn},
\texttt{matplotlib}, \texttt{plotly}, \texttt{beautifulsoup4} where
appropriate.

\subsubsection{Analytic workflow}\label{analytic-workflow-3}

\begin{enumerate}
\def\labelenumi{\arabic{enumi}.}
\tightlist
\item
  Define the document corpus.
\item
  Verify document access and use rights.
\item
  Extract and clean text.
\item
  Build a transition-risk keyword dictionary.
\item
  Apply tokenisation and frequency analysis.
\item
  Use TF-IDF or simple classification where appropriate.
\item
  Compare transition-risk signals across documents.
\item
  Produce a document-risk summary and interpretation.
\end{enumerate}

\subsubsection{Expected outputs}\label{expected-outputs-3}

\begin{itemize}
\tightlist
\item
  Cleaned text corpus.
\item
  Transition-risk dictionary.
\item
  Keyword and TF-IDF analysis.
\item
  Document comparison table.
\item
  Quarto text analytics report.
\end{itemize}

\subsubsection{Dashboard / report
expectations}\label{dashboard-report-expectations-3}

The report should include corpus description, text-cleaning method,
transition-risk framework, key findings, examples of document-level
signals, limitations and implications for investors.

\subsubsection{Assessment criteria}\label{assessment-criteria-3}

\begin{itemize}
\tightlist
\item
  Appropriateness of corpus selection.
\item
  Quality of text preprocessing.
\item
  Transparency of keyword dictionary.
\item
  Validity of interpretation.
\item
  Careful handling of automated text-analysis limitations.
\item
  Professional communication.
\end{itemize}

\subsubsection{Extension potential}\label{extension-potential-3}

The project can support ESG research, policy analysis, transition-risk
screening, investor reporting, journal publication or development of a
larger document-intelligence tool.

\section{Recommended Learning
Resources}\label{recommended-learning-resources}

The resource list should be finalised by the instructor and adapted to
country context, institutional data access, commercial permissions and
the learner's project focus. The following categories and widely known
resources are recommended.

{\def\LTcaptype{none} % do not increment counter
\begin{longtable}[]{@{}
  >{\raggedright\arraybackslash}p{(\linewidth - 4\tabcolsep) * \real{0.3333}}
  >{\raggedright\arraybackslash}p{(\linewidth - 4\tabcolsep) * \real{0.3333}}
  >{\raggedright\arraybackslash}p{(\linewidth - 4\tabcolsep) * \real{0.3333}}@{}}
\toprule\noalign{}
\begin{minipage}[b]{\linewidth}\raggedright
Resource Category
\end{minipage} & \begin{minipage}[b]{\linewidth}\raggedright
Recommended Resource Type
\end{minipage} & \begin{minipage}[b]{\linewidth}\raggedright
Use
\end{minipage} \\
\midrule\noalign{}
\endhead
\bottomrule\noalign{}
\endlastfoot
Python foundations & Python official documentation and tutorial &
Syntax, functions, control flow and standard library. \\
Jupyter & JupyterLab documentation & Notebook workflow and interactive
analysis. \\
pandas & pandas documentation and user guide & Data cleaning,
transformation, joins and time-series handling. \\
NumPy & NumPy documentation & Numerical operations and arrays. \\
Visualisation & Matplotlib, seaborn and Plotly documentation & Static,
statistical and interactive visualisation. \\
Statistics and econometrics & statsmodels documentation & Regression,
diagnostics and time-series modelling. \\
Machine learning & scikit-learn documentation & Supervised learning,
preprocessing, pipelines and evaluation. \\
Geospatial analysis & GeoPandas, Shapely and Folium documentation &
Spatial joins, mapping and interactive maps. \\
Climate data & National climate portals, environmental agencies and open
climate-data platforms & Physical risk, flood risk and exposure
indicators, to be verified. \\
Finance data & Central banks, statistical offices, stock exchanges,
finance-data APIs or commercial data providers & REITs, interest rates,
inflation and market data. \\
Real estate data & Land registry, property portals, valuation datasets,
institutional databases or commercial providers & Transaction and
valuation analytics, subject to rights. \\
ESG and transition-risk documents & Company reports, regulator
documents, planning documents and policy portals & NLP and
transition-risk signal analysis. \\
Dashboards & Streamlit or Dash documentation & Applied dashboard
development. \\
Reproducible reporting & Quarto documentation & Technical reports,
project documentation and publication-ready outputs. \\
Version control & GitHub Docs and Git documentation & Repository
management, documentation and collaboration. \\
\end{longtable}
}

Resource verification requirements:

\begin{itemize}
\tightlist
\item
  Dataset access must be checked before teaching begins.
\item
  Commercial data must not be redistributed without permission.
\item
  Sensitive property-level data must be anonymised or replaced with
  synthetic examples.
\item
  Public datasets must be checked for licensing and citation
  requirements.
\item
  Readings should be updated to include current real estate climate-risk
  literature, flood-risk capitalisation studies, transition-risk
  scholarship and relevant professional standards.
\end{itemize}

\section{Reproducible Workflow and MLOps
Introduction}\label{reproducible-workflow-and-mlops-introduction}

This curriculum introduces reproducible workflow and beginner-level
MLOps as practical habits rather than advanced engineering abstractions.

\subsection{Folder Structure}\label{folder-structure}

The learner should use a consistent folder structure for all projects.
Raw data should be separated from processed data. Scripts, notebooks,
reports, figures and dashboard files should be clearly organised. Each
project should include a \texttt{README.md} explaining the problem, data
sources, workflow, dependencies and outputs.

\subsection{Notebook Discipline}\label{notebook-discipline}

Notebooks should be used for exploration and teaching, but they must
remain readable and reproducible. Each notebook should have a title,
purpose statement, data source note, cleaning section, analysis section,
interpretation section and limitations section. Temporary experiments
should be separated from final analytical notebooks.

\subsection{Quarto Reports}\label{quarto-reports}

Quarto should be used to produce professional reports that integrate
code, tables, figures, interpretation and methodological notes. By the
end of the programme, the learner should be able to prepare HTML and PDF
reports suitable for academic, consultancy or internal professional
audiences.

\subsection{Git and GitHub}\label{git-and-github}

Git should be introduced as a way of tracking work, not as an abstract
software-development tool. The learner should practise committing
meaningful changes, writing commit messages, pushing to GitHub,
maintaining a README and using repositories as a professional portfolio.

\subsection{Documentation}\label{documentation}

Documentation should include data dictionaries, cleaning logs, modelling
assumptions, package requirements, project limitations and stakeholder
interpretation notes. Documentation is assessed because it determines
whether the work can support research, consultancy or product
development.

\subsection{Environment Management}\label{environment-management}

The learner should record package dependencies using
\texttt{requirements.txt}, \texttt{environment.yml} or an equivalent
environment file. This supports reproducibility and reduces the risk
that a project cannot be rerun later.

\subsection{Model Versioning Concepts}\label{model-versioning-concepts}

At an introductory level, the learner should understand that models have
versions, training data, assumptions, parameters, evaluation metrics and
intended-use boundaries. A simple model card should be used for Week 7
and the final capstone where relevant.

\subsection{Basic Deployment Concepts}\label{basic-deployment-concepts}

Deployment should be introduced through Streamlit or Dash. The learner
should understand the difference between a local prototype, a shared
internal dashboard and a public-facing deployed tool. Public deployment
should only occur where data rights, confidentiality, security and
ethical considerations have been addressed.

\subsection{Ethical Data Handling}\label{ethical-data-handling}

Property-level, financial and location-based data may raise privacy,
commercial and ethical issues. The learner should avoid exposing
sensitive records, should anonymise where needed and should document
data provenance and access rights.

\subsection{Reproducible Research
Practices}\label{reproducible-research-practices}

The programme should encourage transparent workflows, documented
assumptions, clear code, version-controlled projects, reproducible
reports and cautious interpretation. These practices are essential for
journal publications, grant-funded projects, consultancy and
professional analytics.

\section{Professional Outputs and
Portfolio}\label{professional-outputs-and-portfolio}

By the end of the programme, the learner should have the following
portfolio outputs:

{\def\LTcaptype{none} % do not increment counter
\begin{longtable}[]{@{}
  >{\raggedright\arraybackslash}p{(\linewidth - 4\tabcolsep) * \real{0.3333}}
  >{\raggedright\arraybackslash}p{(\linewidth - 4\tabcolsep) * \real{0.3333}}
  >{\raggedright\arraybackslash}p{(\linewidth - 4\tabcolsep) * \real{0.3333}}@{}}
\toprule\noalign{}
\begin{minipage}[b]{\linewidth}\raggedright
Output
\end{minipage} & \begin{minipage}[b]{\linewidth}\raggedright
Description
\end{minipage} & \begin{minipage}[b]{\linewidth}\raggedright
Professional Value
\end{minipage} \\
\midrule\noalign{}
\endhead
\bottomrule\noalign{}
\endlastfoot
Cleaned datasets & Analysis-ready property, risk, financial, spatial or
text datasets. & Evidence of data-management competence. \\
Data-quality report & Missing values, duplicates, invalid values, joins
and cleaning decisions. & Supports auditability and professional
reliability. \\
EDA notebook & Exploratory analysis with charts, summaries and
interpretation. & Demonstrates ability to diagnose data patterns. \\
Regression report & Quarto report estimating a pricing or
risk-capitalisation model. & Supports research and valuation
analysis. \\
Machine-learning notebook & Predictive or classification model with
evaluation and model card. & Demonstrates applied predictive
analytics. \\
Geospatial flood-risk map & Static or interactive map showing exposure
patterns. & Supports spatial risk communication. \\
NLP transition-risk analysis & Text analytics of ESG, policy or planning
documents. & Supports transition-risk research and disclosure
analysis. \\
Dashboard prototype & Streamlit or Dash decision-support interface. &
Supports consultancy, stakeholder communication and product
development. \\
GitHub repository & Organised, documented and version-controlled
portfolio. & Demonstrates reproducible professional workflow. \\
Quarto technical report & Integrated analysis report with code, tables,
figures and interpretation. & Suitable for academic, consultancy or
grant outputs. \\
Capstone presentation & Professional presentation of final project. &
Supports stakeholder communication and project defence. \\
\end{longtable}
}

\section{Progression Route Beyond 12
Weeks}\label{progression-route-beyond-12-weeks}

The 12-week curriculum establishes a beginner-to-intermediate applied
Python foundation. The learner's next pathway should move toward
advanced research and professional analytics.

{\def\LTcaptype{none} % do not increment counter
\begin{longtable}[]{@{}
  >{\raggedright\arraybackslash}p{(\linewidth - 4\tabcolsep) * \real{0.3333}}
  >{\raggedright\arraybackslash}p{(\linewidth - 4\tabcolsep) * \real{0.3333}}
  >{\raggedright\arraybackslash}p{(\linewidth - 4\tabcolsep) * \real{0.3333}}@{}}
\toprule\noalign{}
\begin{minipage}[b]{\linewidth}\raggedright
Advanced Area
\end{minipage} & \begin{minipage}[b]{\linewidth}\raggedright
Purpose
\end{minipage} & \begin{minipage}[b]{\linewidth}\raggedright
Possible Learning Direction
\end{minipage} \\
\midrule\noalign{}
\endhead
\bottomrule\noalign{}
\endlastfoot
Advanced econometrics & Strengthen causal and empirical modelling for
real estate research. & Panel data, fixed effects, instrumental
variables and robust inference. \\
Causal inference & Improve interpretation of climate-risk and policy
effects. & Difference-in-differences, matching, synthetic control and
event studies. \\
Spatial econometrics & Address spatial dependence in property markets
and climate exposure. & Spatial lag models, spatial error models and
geographically weighted regression. \\
Bayesian modelling & Model uncertainty and probabilistic risk. &
Bayesian regression, hierarchical models and uncertainty intervals. \\
Advanced machine learning & Improve predictive modelling and feature
engineering. & Gradient boosting, model explainability and pipeline
optimisation. \\
Deep learning & Use only where justified by data and problem scale. &
Image, satellite, remote-sensing or large-scale text applications. \\
Climate scenario analysis & Link assets to future risk pathways. &
Climate scenarios, physical-risk projections and transition pathways. \\
Asset-level risk modelling & Build granular property and portfolio risk
tools. & Asset scoring, resilience indicators and valuation adjustment
frameworks. \\
Portfolio risk analytics & Extend from asset-level to portfolio-level
decision-making. & Diversification, exposure concentration and scenario
stress testing. \\
Dashboard and product development & Move from prototype to professional
tool. & User testing, deployment, authentication and cloud hosting. \\
Cloud deployment & Share tools securely and scalably. & Streamlit
Community Cloud, cloud platforms or institutional servers, subject to
governance. \\
Research publication pipeline & Convert analytics into scholarly
outputs. & Manuscript structure, methods reporting, reproducible
appendices and data documentation. \\
\end{longtable}
}

\section{Alignment Matrix}\label{alignment-matrix}

{\def\LTcaptype{none} % do not increment counter
\begin{longtable}[]{@{}
  >{\raggedright\arraybackslash}p{(\linewidth - 6\tabcolsep) * \real{0.2500}}
  >{\raggedright\arraybackslash}p{(\linewidth - 6\tabcolsep) * \real{0.2500}}
  >{\raggedright\arraybackslash}p{(\linewidth - 6\tabcolsep) * \real{0.2500}}
  >{\raggedright\arraybackslash}p{(\linewidth - 6\tabcolsep) * \real{0.2500}}@{}}
\toprule\noalign{}
\begin{minipage}[b]{\linewidth}\raggedright
Curriculum Requirement
\end{minipage} & \begin{minipage}[b]{\linewidth}\raggedright
Curriculum Response
\end{minipage} & \begin{minipage}[b]{\linewidth}\raggedright
Where It Appears
\end{minipage} & \begin{minipage}[b]{\linewidth}\raggedright
Expected Learner Outcome
\end{minipage} \\
\midrule\noalign{}
\endhead
\bottomrule\noalign{}
\endlastfoot
Non-generic Python training aligned with real estate and climate-risk
analytics & Domain-specific examples are embedded from Week 1 onward. &
Programme rationale, curriculum map, weekly plans & Learner studies
Python through relevant professional problems. \\
Target learner is a researcher, academic and real estate
finance/climate-risk professional & Learner profile and pedagogy are
designed for advanced conceptual expertise and beginner programming. &
Learner profile, design principles, instructor notes & Learner receives
appropriate pacing and professional relevance. \\
Beginner-to-intermediate Python with advanced pathway & Foundations are
taught first, followed by applied analytics and post-programme
progression. & Weeks 1-12, progression route & Learner develops from
basic syntax to applied analytics readiness. \\
Strong quantitative and statistical background & Regression,
econometrics, time series and machine learning are included without
over-explaining basic research logic. & Weeks 6-8, capstones & Learner
converts statistical knowledge into Python workflows. \\
Need to clean and manage large property, climate, transaction and
financial datasets & Data cleaning, pandas, NumPy and quality assurance
are central. & Weeks 3-4 & Learner can prepare analysis-ready
datasets. \\
Excel, CSV, APIs, geospatial files and public datasets & Data import and
data-source strategy cover these formats. & Week 3, dataset strategy &
Learner can work across common data sources. \\
Exploratory data analysis & EDA is explicitly taught with professional
visual outputs. & Week 5 & Learner can diagnose and communicate data
patterns. \\
Professional charts, tables and dashboards & Visualisation and dashboard
development are core components. & Weeks 5 and 11 & Learner can produce
stakeholder-ready outputs. \\
Statistical and econometric modelling & Regression and real estate
pricing models are included. & Week 6 & Learner can estimate and
interpret pricing/risk models. \\
Machine learning for exposure, discounts, vulnerability and resilience &
ML week focuses on prediction, classification and risk scoring. & Week 7
& Learner can build and evaluate applied predictive models. \\
Time-series analysis for REITs, housing prices, rates, inflation and
shocks & Time-series module includes returns, volatility and event
indicators. & Week 8 & Learner can analyse temporal market and shock
dynamics. \\
Geospatial analysis & Spatial joins, mapping and flood-exposure analysis
are included. & Week 9 & Learner can map and analyse property
exposure. \\
Climate-risk and flood-risk data integration & Property, risk, local
authority and vulnerability data are joined and modelled. & Weeks 3, 4,
6, 9, 11 & Learner can integrate multi-source risk evidence. \\
NLP for policy, planning, ESG and sentiment & NLP module targets
transition-risk document analysis. & Week 10 & Learner can extract
transition-risk signals from text. \\
Dashboard development with Streamlit or Dash & Dashboard prototype is a
required professional output. & Week 11 & Learner can build a
climate-risk communication tool. \\
MLOps and reproducible workflows & GitHub, Quarto, folder structures,
documentation, environment management and deployment concepts are
included. & Weeks 1, 7, 11, 12, reproducible workflow section & Learner
can produce reproducible research and professional workflows. \\
Practical assessments preferred over theoretical tests & Assessment
design prioritises exercises, projects, portfolio and presentation. &
Assessment design and rubric & Learner is assessed through authentic
outputs. \\
Capstone project options & Four capstone briefs are developed from the
uploaded use cases. & Capstone section & Learner can select a project
aligned with research and professional goals. \\
Detailed syllabus and learning pathway & Full weekly map, teaching
plans, exercises, datasets, assessments and progression route are
included. & Entire document & Learner and instructor have an
implementation-ready curriculum. \\
Future research, consultancy, grant and product development & Portfolio
outputs and extension pathways are included. & Professional outputs,
progression route, capstones & Learner can extend training into academic
and professional outputs. \\
\end{longtable}
}

\section{Implementation Notes for the
Instructor}\label{implementation-notes-for-the-instructor}

This curriculum should be taught as a professional analytics programme,
not as a coding bootcamp. The learner is not a general beginner in
research or quantitative reasoning. The main teaching challenge is to
scaffold programming without oversimplifying the domain.

Pacing should be firm but supportive. Weeks 1-3 must establish
confidence in the environment, syntax and data cleaning. If these weeks
are rushed, later modelling and dashboard work will become fragile.
However, examples should remain domain-specific from the beginning so
that the learner sees immediate relevance.

The instructor should avoid generic datasets unless they are used only
for a very short technical demonstration. Exercises should use property,
valuation, flood-risk, climate-risk, financial, geospatial, ESG or
policy examples. Where real data cannot be used, synthetic data should
be constructed to mirror the structure of real professional datasets.

Debugging should be treated as a normal part of analytic practice. The
learner should be taught to read error messages, inspect objects, check
data types, verify joins and test code incrementally.

Data limitations must be discussed explicitly. Real estate and
climate-risk datasets are often incomplete, proprietary, spatially
inconsistent or governed by licensing restrictions. The learner should
be trained to document limitations rather than hide them.

Syntax should be balanced with interpretation. Every coding exercise
should end with a short professional interpretation: what does the
output mean for valuation, investment, risk, policy, resilience,
insurance, mortgageability or market liquidity?

The instructor should require documentation from the beginning. README
files, data dictionaries, cleaning logs and Quarto notes are not
optional extras. They are part of the evidence standard required for
research, consultancy and professional analytics.

The final capstone should be selected by Week 5 so that exercises from
Weeks 6-11 can gradually feed into the final project. The capstone
should not be treated as a separate end-of-course assignment but as an
integrated analytical system built over time.

\section{Quality-Control Checklist}\label{quality-control-checklist}

Use this checklist before finalising or delivering the programme.

{\def\LTcaptype{none} % do not increment counter
\begin{longtable}[]{@{}
  >{\raggedright\arraybackslash}p{(\linewidth - 2\tabcolsep) * \real{0.5000}}
  >{\raggedright\arraybackslash}p{(\linewidth - 2\tabcolsep) * \real{0.5000}}@{}}
\toprule\noalign{}
\begin{minipage}[b]{\linewidth}\raggedright
Requirement
\end{minipage} & \begin{minipage}[b]{\linewidth}\raggedright
Confirmed
\end{minipage} \\
\midrule\noalign{}
\endhead
\bottomrule\noalign{}
\endlastfoot
The programme is structured as 12 weeks / 90 days. & Yes \\
The weekly workload assumes approximately 4-6 hours per week. & Yes \\
The curriculum begins with Python foundations. & Yes \\
The curriculum progresses into applied analytics. & Yes \\
The curriculum is designed for a PhD-trained real estate finance and
climate-risk professional. & Yes \\
The curriculum avoids generic Python training. & Yes \\
Real estate and climate-risk use cases are embedded throughout. & Yes \\
Flood risk is included. & Yes \\
Physical climate risk is included. & Yes \\
Transition risk is included. & Yes \\
Real estate pricing and valuation are included. & Yes \\
REITs and financial-market shocks are included. & Yes \\
Infrastructure shocks and resilience are included. & Yes \\
Investment decision-making is included. & Yes \\
Python foundations include variables, data types, functions, loops and
conditionals. & Yes \\
Data analysis with pandas and NumPy is included. & Yes \\
Visualisation with matplotlib, seaborn and plotly is included. & Yes \\
Excel, CSV, API and web-sourced data workflows are included. & Yes \\
Statistical analysis and regression modelling are included. & Yes \\
Machine learning for prediction and classification is included. & Yes \\
Time-series analysis is included. & Yes \\
Geospatial analysis is included. & Yes \\
NLP for ESG, policy, planning and transition-risk documents is included.
& Yes \\
Dashboard development with Streamlit or Dash is included. & Yes \\
GitHub, notebooks, documentation, Quarto and reproducible workflow are
included. & Yes \\
Assessment tasks are practical rather than purely theoretical. & Yes \\
Weekly coding exercises are included. & Yes \\
Short applied tasks are included. & Yes \\
A mid-programme mini-project is included. & Yes \\
A final capstone project is included. & Yes \\
Four capstone options are provided. & Yes \\
Dataset types are identified without fabricating access. & Yes \\
Recommended resources are provided as verifiable categories and
documentation sources. & Yes \\
Professional outputs and portfolio artifacts are specified. & Yes \\
A progression route beyond 12 weeks is included. & Yes \\
Instructor implementation notes are included. & Yes \\
The curriculum is packaged as a Quarto \texttt{.qmd} document. & Yes \\
\end{longtable}
}




\end{document}
